\documentclass[a4paper,11pt]{article}
        \usepackage[top=15mm,bottom=18mm,left=16mm,right=16mm]{geometry}
        \usepackage{fontspec}
        \usepackage{xeCJK}
        \setmainfont{Times New Roman}
        \setCJKmainfont{Yu Gothic}
        \setmonofont{Consolas}
        \setCJKmonofont{Yu Gothic}
        \usepackage{booktabs}
        \usepackage{longtable}
        \usepackage{tabularx}
        \usepackage{array}
        \usepackage{enumitem}
        \usepackage{hyperref}
        \usepackage{listings}
        \usepackage{xcolor}
        \hypersetup{
          colorlinks=true,
          linkcolor=blue,
          urlcolor=blue,
          pdftitle={Windows入口 PowerShell ルータ},
          pdfauthor={Codex}
        }
        \lstset{
          basicstyle=\ttfamily\small,
          breaklines=true,
          columns=fullflexible,
          keepspaces=true,
          frame=single,
          framerule=0.2pt,
          rulecolor=\color{black},
          showstringspaces=false
        }
        \setlength{\parskip}{0.42em}
        \setlength{\parindent}{1em}
        \renewcommand{\arraystretch}{1.15}
        \title{01. Windows入口 PowerShell ルータ}
        \author{solar\_ws0129-main}
        \date{2026-07-29}
        \begin{document}
        \maketitle

        \section*{対象}
        \begin{tabularx}{\linewidth}{>{\raggedright\arraybackslash}p{0.18\linewidth} X}
        \toprule
        項目 & 内容 \\
        \midrule
        ファイル & \path|SolarSim.ps1| \\
        種別 & PowerShell \\
        区分 & 入口 \\
        役割 & Windows PowerShell から build、launch、offline script、Grafana 操作をまとめて呼び出す最上位入口。 \\
        起動文脈 & オペレータが最初に叩くコマンド入口。 \\
        \bottomrule
        \end{tabularx}

        \section*{このファイルを読む理由}
        Windows PowerShell から build、launch、offline script、Grafana 操作をまとめて呼び出す最上位入口。

        \begin{itemize}[leftmargin=1.6em]
        \item WSL ディストリビューションを解決する。
\item Action と Mode を shell 側へ受け渡す。
\item dashboard や report の自動オープンも担う。
        \end{itemize}

        \section*{呼び出し関係}
        \subsection*{どこから呼ばれるか}
        \begin{itemize}[leftmargin=1.6em]
        \item 人間の運用操作
        \end{itemize}

        \subsection*{次に読むと理解が進むファイル}
        \begin{itemize}[leftmargin=1.6em]
        \item \path|scripts/solar_control.sh|
\item \path|launch/solar_race_live_wifi.launch.py|
\item \path|scripts/solar_sim.py|
        \end{itemize}

        \subsection*{同一リポジトリ内の主要依存}
        \begin{itemize}[leftmargin=1.6em]
        \item 同一リポジトリ内の明示 import は少ない。
        \end{itemize}

        \section*{ソース冒頭の把握}
        モジュール docstring または先頭意図の要約:

        モジュール docstring は明示されていない。

        \subsection*{代表コード断片}
        \begin{lstlisting}
param(
    [ValidateSet('up', 'build', 'start', 'stop', 'restart', 'status', 'graph', 'simulate', 'historical-weather', 'historical-simulate', 'forecast', 'identify', 'fit', 'learn', 'audit', 'package', 'blank', 'log', 'grafana', 'grafana-stop')]
    [string]$Action = 'up',
    [ValidateSet('sim', 'measure', 'live', 'live_wifi')]
    [string]$Mode = 'sim',
    [string]$Profile = 'project_packages/bwsc2025_fitted_mle19_energywindow_inertia/profile.yaml',
    [string]$Distro = '',
    [switch]$Help
)

$ErrorActionPreference = 'Stop'

if ($Help) {
    @'
SolarSim.ps1 -Action <action> [-Mode <mode>] [-Profile <profile.yaml>] [-Distro <WSL name>]

Actions: up, build, start, stop, restart, status, graph, simulate, forecast,
         identify, fit, learn, audit, package, blank, log, grafana, grafana-stop
Modes:   sim, measure, live, live_wifi

Examples:
  .\SolarSim.ps1 -Action build
  .\SolarSim.ps1 -Action up -Mode sim
  .\SolarSim.ps1 -Action simulate -Profile project_packages/<vehicle>/profile.yaml
'@ | Write-Host
    exit 0
}

function Normalize-WslText {
    param([string]$Text)
    if ($null -eq $Text) {
        return ''
    }
    return ([string]$Text).Replace("`0", '').Trim()
}

function Resolve-WslDistro {
    param([string]$Preferred)
    if ($Preferred) {
        return (Normalize-WslText $Preferred)
    }
    $installed = (& wsl.exe -l -q) | ForEach-Object { Normalize-WslText $_ } | Where-Object { $_ }
    foreach ($candidate in @('Ubuntu-22.04', 'Ubuntu', 'Ubuntu-24.04')) {
        if ($installed -contains $candidate) {
            return $candidate
        }
    }
    if ($installed.Count -gt 0) {
        return $installed[0]
    }
    throw 'WSL distro was not found.'
}

function Resolve-WslArgPath {
    param(
        [string]$TargetDistro,
        [string]$PathValue
    )
    if ([string]::IsNullOrWhiteSpace($PathValue)) {
        return ''
    }
    if ([System.IO.Path]::IsPathRooted($PathValue)) {
        return (& wsl.exe -d $TargetDistro wslpath -a $PathValue).Trim()
    }
    return ($PathValue -replace '\\', '/')
}

function Invoke-WslRepoCommand {
    param(
        [string]$RepoRoot,
\end{lstlisting}


        \section*{主要構造の読み取り}
        action 分岐は 16 件。



        \subsection*{ファイルを上から読んだときの定義順}
            \begin{longtable}{p{0.07\linewidth} p{0.84\linewidth}}
            \toprule
            行 & 内容 \\
            \midrule
            1 & param( \\
2 & [ValidateSet('up', 'build', 'start', 'stop', 'restart', 'status', 'graph', 'simulate', 'historical-weather', 'historical-simulate', 'forecast', 'identify', 'fit', 'learn', 'audit', 'package', 'blank', 'log', 'grafana', 'grafana-stop')] \\
3 & [string]\$Action = 'up', \\
4 & [ValidateSet('sim', 'measure', 'live', 'live\_wifi')] \\
5 & [string]\$Mode = 'sim', \\
6 & [string]\$Profile = 'project\_packages/bwsc2025\_fitted\_mle19\_energywindow\_inertia/profile.yaml', \\
7 & [string]\$Distro = '', \\
8 & [switch]\$Help \\
9 & ) \\
11 & \$ErrorActionPreference = 'Stop' \\
13 & if (\$Help) \{ \\
14 & @' \\
15 & SolarSim.ps1 -Action <action> [-Mode <mode>] [-Profile <profile.yaml>] [-Distro <WSL name>] \\
17 & Actions: up, build, start, stop, restart, status, graph, simulate, forecast, \\
18 & identify, fit, learn, audit, package, blank, log, grafana, grafana-stop \\
19 & Modes:   sim, measure, live, live\_wifi \\
21 & Examples: \\
22 & .\textbackslash{}SolarSim.ps1 -Action build \\
23 & .\textbackslash{}SolarSim.ps1 -Action up -Mode sim \\
24 & .\textbackslash{}SolarSim.ps1 -Action simulate -Profile project\_packages/<vehicle>/profile.yaml \\
25 & '@ | Write-Host \\
26 & exit 0 \\
27 & \} \\
29 & function Normalize-WslText \{ \\
30 & param([string]\$Text) \\
31 & if (\$null -eq \$Text) \{ \\
32 & return '' \\
33 & \} \\
34 & return ([string]\$Text).Replace("`0", '').Trim() \\
35 & \} \\
37 & function Resolve-WslDistro \{ \\
38 & param([string]\$Preferred) \\
39 & if (\$Preferred) \{ \\
40 & return (Normalize-WslText \$Preferred) \\
41 & \} \\
42 & \$installed = (\& wsl.exe -l -q) | ForEach-Object \{ Normalize-WslText \$\_ \} | Where-Object \{ \$\_ \} \\
43 & foreach (\$candidate in @('Ubuntu-22.04', 'Ubuntu', 'Ubuntu-24.04')) \{ \\
44 & if (\$installed -contains \$candidate) \{ \\
45 & return \$candidate \\
46 & \} \\
            \bottomrule
            \end{longtable}



        \section*{関数・クラスを上から順に解説}
        このファイルでは独立した関数・クラス定義が少ない。









        \subsection*{shell 分岐と外部コマンド}
            \begin{longtable}{p{0.07\linewidth} p{0.18\linewidth} p{0.67\linewidth}}
            \toprule
            行 & 種別 & 内容 \\
            \midrule
            250 & action & \path|up| \\
278 & action & \path|build| \\
281 & action & \path|start| \\
284 & action & \path|stop| \\
287 & action & \path|restart| \\
290 & action & \path|status| \\
293 & action & \path|graph| \\
299 & action & \path|simulate| \\
306 & action & \path|historical-weather| \\
309 & action & \path|historical-simulate| \\
317 & action & \path|forecast| \\
320 & action & \path|identify| \\
323 & action & \path|fit| \\
326 & action & \path|learn| \\
329 & action & \path|audit| \\
332 & action & \path|log| \\
42 & command & \path!$installed = (& wsl.exe -l -q) | ForEach-Object { Normalize-WslText $_ } | Where-Object { $_ }! \\
63 & command & \path|return (& wsl.exe -d $TargetDistro wslpath -a $PathValue).Trim()| \\
76 & command & \path|$wslRepo = (& wsl.exe -d $TargetDistro wslpath -a $RepoRoot).Trim()| \\
81 & command & \path|& wsl.exe -d $TargetDistro bash -lc "cd '$wslRepo' && bash './scripts/solar_control.sh' '$Command' '$TargetMode' '$wslProfile'"| \\
207 & command & \path|if (-not (Get-Command docker -ErrorAction SilentlyContinue)) {| \\
214 & command & \path|& docker compose up -d| \\
216 & command & \path|throw "docker compose up failed with exit code $LASTEXITCODE"| \\
220 & command & \path|& docker compose down| \\
222 & command & \path|throw "docker compose down failed with exit code $LASTEXITCODE"| \\
235 & command & \path|& python $builder --force| \\
257 & command & \path|if (Get-Command docker -ErrorAction SilentlyContinue) {| \\
260 & command & \path|& docker compose up -d| \\
262 & command & \path|throw "docker compose up failed with exit code $LASTEXITCODE"| \\
            \bottomrule
            \end{longtable}

        \section*{処理の流れ}
        \begin{enumerate}[leftmargin=1.7em]
        \item PowerShell 引数を解釈する。
\item WSL 側のパスへ変換する。
\item scripts/solar\_control.sh を action と mode 付きで呼ぶ。
\item 必要に応じて dashboard、graph、report を開く。
        \end{enumerate}

        \section*{このファイルの位置づけ}
        この資料群では、\path|SolarSim.ps1| を
        \textbf{入口}の中核として扱う。
        したがって、ここで扱う処理は単独で完結せず、
        前段の profile / launch / telemetry と後段の planner / logger / report へ繋がる。

        \end{document}
